% Source-aligned English for printed p.191.
\clearpage
\SourcePage{191}

\LayerTarget{tohoku:section:4.3}
\textbf{4.3. Case of a constant sheaf of rings.}
\LayerTarget{tohoku:setup:constant-ring-sheaf}
Suppose that $\mathbf O$ is the constant sheaf of rings defined by a ring
with identity $O$. Let $M$ be an $O$-module, and let $\mathbf M$ be the
constant $\mathbf O$-module it defines. For every $\mathbf O$-module $A$,
there is then a natural isomorphism
\LayerTarget{tohoku:eq:4.3.1}
\[
\tag{4.3.1}
\mathrm{Hom}_{\mathbf O}(\mathbf M,A)=\mathrm{Hom}_O(M,\Gamma(A)).
\]
This is a functorial isomorphism, which may also be written
\LayerTarget{tohoku:eq:4.3.2}
\[
\tag{4.3.2}
h_{\mathbf M}=h_M\Gamma
\]
\LayerTarget{tohoku:setup:functors-h-constant-module}
where, as in the preceding subsection, we set
$h_{\mathbf M}(A)=\mathrm{Hom}_{\mathbf O}(\mathbf M,A)$ and
$h_M(N)=\mathrm{Hom}_O(M,N)$ for every $\mathbf O$-module $A$ and every
$O$-module $N$. Here $\Gamma$ is regarded as a functor from
$\mathbf C^{\mathbf O}$ to the category $\mathbf C^O$ of left $O$-modules, and
$h_M$ as a functor from $\mathbf C^O$ to the category of abelian
groups. Theorem 2.4.1 applies, by virtue of the following

\LayerTarget{tohoku:lemma:4.3:gamma-injective}
\textsc{Lemma} \textit{If $A$ is an injective $\mathbf O$-module,
then $\Gamma(A)$ is an injective $O$-module.}

\LayerTarget{tohoku:proof:lemma-4.3-gamma-injective}
We must show that the functor
$M\to\mathrm{Hom}_O(M,\Gamma(A))$ is exact. This follows at once from
(4.3.1), since $M\to\mathbf M$ is an exact functor. Thus, by Theorem
2.4.1, and observing that the derived functors of $\Gamma$ are the same
whether $\Gamma$ is regarded as taking its values in $\mathbf C^O$ or in the category
$\mathbf C$ of abelian groups, we obtain:

\LayerTarget{tohoku:theorem:4.3.1}
\textsc{Theorem 4.3.1.} \textit{Let $X$ be a space, $O$ a ring with
identity, and $M$ an $O$-module; let $\mathbf O$ and $\mathbf M$ be the
constant sheaves on $X$ defined by $O$ and $M$. There is then a cohomological
spectral functor on $\mathbf C^{\mathbf O}$, abutting to the graded functor
$(\mathrm{Ext}_{\mathbf O}^p(X;\mathbf M,A))$, whose initial term is}
\LayerTarget{tohoku:eq:4.3.3}
\[
\tag{4.3.3}
\mathrm{II}_2^{p,q}(A)=\mathrm{Ext}_O^q(M,H^p(X,A)).
\]

\LayerTarget{tohoku:remark:edge-maps-and-five-term-sequence-4.3.1}
As usual, this yields ``edge-homomorphisms'' and a five-term exact sequence, which
the reader may write out.
\LayerTarget{tohoku:construction:computing-constant-sheaf-spectral-functor}
To compute the spectral functor, we shall use Corollary 1 to Proposition
2.5.3, employing the resolvent functor $\Gamma \mathrm C(A)$ for $\Gamma$, regarded
here as a covariant functor from $\mathbf C^{\mathbf O}$ to $\mathbf C^O$, where $\mathrm C(A)$ is the
canonical resolution of $A$ (cf. 2.3), and the resolvent functor $h_L$ for
$h_M$, where $L$ is a left resolution of $M$ by free $O$-modules. We obtain
the first part of the following

\LayerTarget{tohoku:proposition:4.3.2}
\textsc{Proposition 4.3.2.} \textit{Under the hypotheses of Theorem
\textup{4.3.1}, choose a left resolution $L$ of $M$ by free $O$-modules. Then the
spectral functor of Theorem \textup{4.3.1} is computed by taking the second spectral
sequence of the bicomplex $\mathrm{Hom}_O(L,\Gamma \mathrm C(A))$, where $\mathrm C(A)$ is
the canonical resolution of $A$ \textup{(}cf.\ \textup{2.3)}. If the $L_i$ are isomorphic to
$O^{n_i}$, where $n_i$ is an integer $>0$, the first spectral sequence of this
bicomplex gives the spectral functor of Theorem \textup{4.2.1}.}

\LayerTarget{tohoku:proof:proposition-4.3.2}
For the latter assertion, it suffices to consider the resolution $\mathbf L$
of $\mathbf M$ deduced from $L$ and apply Corollary 3 to Proposition 4.2.4.
\LayerTarget{tohoku:reduction:kunneth-for-constant-sheaves}
The preceding proposition reduces the computation of
$\mathrm{Ext}_{\mathbf O}(\mathbf M,A)$ and its spectral sequences to a
classical problem of Künneth type. Thus:

\TohokuContinuation{192}
